\section{Theory and motivation}
\label{sec:lit}
\subsection*{Institutional constraints, incarceration rates and equity concerns}

Incarceration rates and sentencing outcomes are shaped by the interaction of multiple institutional actors. Local crime rates—often endogenous—affect the volume and composition of cases entering the system. Legislatures define sentencing ranges through substantive criminal law, while the executive influences outcomes by setting policing priorities and appointing prosecutors. Judicial actors, however, play a central role through discretionary decisions: whether individuals face pretrial detention, are found guilty or acquitted, receive custodial sentences or fines, and, when applicable, the length of incarceration. 

A large literature shows that these judicial decisions respond systematically to extralegal factors. Persistent differences in judge preferences and characteristics \citep{Kling-Random, Aizer-Juvenile}, political incentives \citep{Gordon-Judges-Election, Berdejo-PoliticalCycles}, institutional constraints such as prison overcrowding \citep{Schnepel}, private incentives \citep{Galinato-2020, Mukherjee-Private, Dippel}, fiscal considerations \citep{Ouss-Incentives}, and broader systemic pressures—including judicial vacancies \citep{Yang-Resource-AEJ}, quasi-random caseload shocks \citep{JudgeCaseloads-JPE-Shumway}, and the organization of law enforcement agencies \citep{Ater-Police-Org-JPE}—all shape sentencing outcomes. Together, these findings underscore that sentencing is not a mechanical application of legal rules, but reflects the institutional environment in which judges operate.

From an equity perspective, this discretion matters because judicial decisions directly shape life trajectories—affecting future employment, social integration, family stability, and access to rights. A growing body of evidence shows that the conditions under which individuals are confined are central to these outcomes. Reduced overcrowding, lower-security environments, and access to rehabilitation programs are associated with lower recidivism and better post-release outcomes \citep{Tobon, Mastro, Chen, Williams}, while diversion away from incarceration improves labor market prospects even when sentence length itself has limited effects \citep{Schnepel, Kling-Random}. Importantly, these equity costs arise not only from final custodial sentences, but also from earlier judicial choices—particularly the use of pretrial detention. Individuals detained prior to trial often experience job loss, income shocks, and family disruption, despite not having been found guilty, and are exposed to prison conditions and peer environments that may have lasting consequences.

These concerns are especially salient in high-violence contexts such as Mexico. Individuals accused of relatively minor offenses may be detained or incarcerated alongside organized crime members, increasing exposure to criminal networks even when cases ultimately result in acquittal or non-custodial sentences. Longer periods of detention or incarceration raise the likelihood of sustained interaction with these groups, potentially reinforcing cycles of criminal involvement after release \citep{Pintoff}. Moreover, rehabilitation services in many Mexican prisons remain limited, so extended exposure to incarceration—whether pretrial or post-conviction—may further reduce the prospects of successful reintegration \citep{Evalua2}.

Against this backdrop, why might local judges respond to the construction of federal prisons near the municipalities where they sentence individuals? One mechanism operates through capacity constraints. During periods of overcrowding, judges may have been more reluctant to impose custodial sentences because distance and overcapacity are perceived as additional punishment. Although judges do not directly bear incarceration costs, they are legally required to assign individuals to the nearest facility, making local capacity salient. Once new capacity becomes available, sentencing decisions may revert toward baseline levels, mechanically increasing incarceration \citep{Schnepel}. 

A second mechanism relates to caseload pressures: higher policing intensity may raise the number of cases without changing sentencing distributions, or alternatively encourage faster resolution through non-custodial outcomes \citep{Yang-Resource-AEJ}. The idea is the following: when dockets are heavy, defense lawyers, prosecutors, and judges all have incentives to arrange plea bargains that reduce sentences or substitute fines to avoid trial (given time constraints). In this context, an increase in policing following a capacity expansion could raise caseloads, but paradoxically result in fewer prison sentences relative to other outcomes, shorter sentences, and more fines. A third and final mechanism is salience and political pressure. Prison construction is highly visible and may generate expectations that new capacity will be used. In Mexico, large federal prisons were politically salient and often framed as responses to insecurity,\footnote{For example, in 2024, a presidential candidate proposed building a massive prison as a response to rising insecurity \citep{FinancieroXochitl}.} potentially translating executive pressure into greater use of pretrial detention or custodial sentences.

To fix ideas, consider the first mechanism with a simple example. Imagine an individual accused of a minor theft offense in a municipality before the construction of a nearby federal prison. Under conditions of severe overcrowding, a judge may be more inclined to release the defendant pending trial or to impose a non-custodial sentence, limiting disruption to employment, family ties, and future opportunities. Now consider the same individual, accused of the same offense with identical observable characteristics, but apprehended after the prison becomes operational. If judicial decisions respond to the availability of new capacity, the judge may be more likely to impose pretrial detention or a custodial sentence, even though the legal merits of the case are unchanged. In the first scenario, the individual may retain employment and social ties; in the second, they may experience job loss, prolonged exposure to prison conditions, and long-term stigma—even if ultimately acquitted or sentenced to a short term. This comparison illustrates how institutional changes such as prison construction can generate meaningful differences in life outcomes through shifts in judicial discretion, rather than through changes in criminal behavior or legal standards. 